\chapter{总结与展望}

\section{工作总结}

本文围绕 2025 Digital Lifestyle Benchmark Dataset 完成了一个完整的大数据分析课程项目。工作内容包括数据集筛选与合规性说明、数据质量检查、工程特征构造、探索性可视化、分类筛查、回归预测、聚类画像和 LaTeX 报告整理。实验结果表明，数字依赖得分在当前特征下具有较强可预测性，高风险标记可以通过偏召回阈值策略进行初筛，而生产力得分预测较弱，聚类结果更适合用于探索性画像。

\section{应用价值}

从课程应用角度看，本文展示了如何在一个主题统一的数据集中同时组织分类、回归和聚类任务。分类部分强调不能只看 Accuracy，需要结合 Recall、F1、PR-AUC 和混淆矩阵选择筛查策略；回归部分说明不同目标变量的可预测性差异；聚类部分展示了如何在不使用 outcome 字段训练的情况下形成用户画像，并用结果变量进行事后解释。

\section{不足与局限性}

本文存在若干局限。首先，数据集为合成数据，不能代表真实人群行为规律。其次，本文使用的是传统机器学习模型和表格特征，没有引入更复杂的时序或上下文信息。第三，分类任务的 PR-AUC 和 F1 仍有提升空间，误报和漏报都存在。第四，聚类 Silhouette=0.1860，说明分群边界较弱，聚类画像只能作为描述性探索。

\section{后续改进}

后续可以从三个方向改进。第一，在真实、合规、匿名化的数据上验证模型稳定性。第二，引入更细粒度的时间特征，例如夜间设备使用、通知集中度和连续睡眠变化。第三，进一步比较模型解释方法和阈值选择策略，使高风险筛查结果更适合人工复核流程。

\section{课程收获与个人反思}

通过本项目，我进一步理解了数据分析不能只追求指标大小，还要关注数据来源、目标泄漏、任务定义和结果边界。特别是在生产力预测为负结果、聚类边界较弱的情况下，诚实报告模型不能支持的结论，比强行写出“显著影响”更重要。这个过程也让我熟悉了将 notebook、src 工具函数、结果 CSV、图表文件和 LaTeX 报告组织为可复现项目的完整流程。
